\documentclass[11pt]{article}
\usepackage{fontspec}
\setmainfont{TeX Gyre Pagella}
\setmonofont{DejaVu Sans Mono}[Scale=0.85]
\usepackage[a4paper,margin=1in]{geometry}
\usepackage{enumitem}
\usepackage{hyperref}
\usepackage{parskip}
\usepackage{amsmath,amssymb}
\usepackage{longtable}
\usepackage{booktabs}
\usepackage{array}
\usepackage{calc}
\usepackage{fancyvrb}
\usepackage{fvextra}
\usepackage{seqsplit}

% pandoc emits \tightlist for compact list items; needs defining since
% we're not using pandoc's own default template.
\providecommand{\tightlist}{%
  \setlength{\itemsep}{0pt}\setlength{\parskip}{0pt}}

% longtable columns produced by pandoc sometimes need slightly relaxed
% column-width arithmetic; this keeps them from overflowing the margin.
\setlength{\LTleft}{0pt}
\setlength{\LTright}{0pt}

\title{\textbf{Buffer World}\\
\large A 3D Spatial Tutorial for Real Vim Motions\\
\normalsize Concept and Technical Specification}
\author{Flyxion\\\small Independent Researcher}
\date{July 2026}

\begin{document}
\maketitle
\thispagestyle{empty}
\newpage

\tableofcontents
\newpage

\part{Concept and Design}

\hypertarget{sec-elevator-pitch}{%
\section{Elevator Pitch}\label{sec-elevator-pitch}}

Vim Adventures (2011) teaches \texttt{hjkl} and friends by turning a 2D tile grid into an
overworld: you move the way you'd move the cursor, and the game only lets you
progress by learning the next motion. This project asks what changes if the
buffer isn't a flat grid but a navigable 3D space --- where ``the text'' is
something you walk through, climb, and reshape, rather than something you look
at from directly above.

The goal stays disciplined: every mechanic maps to a \textbf{real, unmodified Vim
command}. Nothing is vim-\emph{flavored}; it's vim, wearing a world.

\hypertarget{sec-design-pillars}{%
\section{Design Pillars}\label{sec-design-pillars}}

\begin{enumerate}
\def\labelenumi{\arabic{enumi}.}
\tightlist
\item
  \textbf{No invented commands.} If it's not a real Vim motion, operator, or mode,
  it doesn't go in the game. The temptation to add a ``jump boot'' power-up that
  isn't tied to an actual keystroke is the thing to resist hardest.
\item
  \textbf{The 3D-ness has to earn its keep.} A 3D game that plays exactly like the
  2D original with better lighting isn't worth building. Section~\ref{sec-the-core-translation-what-is-3d-space-of-here} below is
  the actual pitch: which Vim concepts \emph{only} make sense once you have a third
  axis.
\item
  \textbf{Mistakes should look wrong, not just fail silently.} In real Vim, a
  mistyped command either does something confusing or nothing at all --- the
  worst possible teacher. The game's job is to make the \emph{consequence} of a
  command visible immediately, in space, before the player has to reason
  about it abstractly.
\item
  \textbf{Keystrokes stay real keystrokes.} No on-screen button mashing. The player
  types \texttt{3dw} on a real keyboard and watches it happen. The game is a skin
  over an actual modal input system, not a simulation of one.
\end{enumerate}

\hypertarget{sec-the-core-translation-what-is-3d-space-of-here}{%
\section{\texorpdfstring{The Core Translation: What Is 3D Space \emph{Of}, Here?}{The Core Translation: What Is 3D Space Of, Here?}}\label{sec-the-core-translation-what-is-3d-space-of-here}}

This is the one decision everything else depends on, so it's worth spelling
out three candidate mappings before picking one.

\textbf{Option A --- Depth as buffer stack.} X/Y is the familiar row/column plane
(rendered as a walkable floor grid, one tile per character), and Z is \emph{which
buffer or split you're in}. Switching windows (\texttt{Ctrl-w} family) becomes
physically walking between elevated platforms; \texttt{:vsplit} cracks the floor open
into a new platform beside you.

\textbf{Option B --- Depth as document structure.} X/Y is still row/column, but Z
encodes nesting --- indentation level, or brace/paren depth. Diving into a
function body means descending a level. \texttt{\%} (jump to matching bracket) becomes
a vertical teleport between the top and bottom of a shaft.

\textbf{Option C --- Depth as time (undo tree).} X/Y is the buffer as usual; Z is the
undo history, rendered as a tower you can walk down (\texttt{u}) and back up
(\texttt{Ctrl-r}), with branches for \texttt{g-}/\texttt{g+} if the undo tree forked.

\textbf{Recommendation: build A first, and treat B and C as world-specific level
gimmicks rather than the permanent camera model.} A is the least conceptually
demanding --- most new Vim users don't yet think about splits, so making splits
\emph{physical} rather than abstract is a genuine teaching win. B and C are
excellent later-level material precisely because they're harder ideas
(nesting, undo trees) that benefit from being made concrete, but building the
whole game around them first would front-load the hardest concepts.

\hypertarget{sec-mode-system-as-world-state}{%
\section{Mode System as World State}\label{sec-mode-system-as-world-state}}

Vim's biggest early confusion for new users isn't any single motion, it's that
\textbf{the same key does different things depending on mode.} This is the single
best argument for going 3D at all: mode can be a \emph{visible property of the
world}, not a mental flag the player has to track.

\begin{itemize}
\tightlist
\item
  \textbf{Normal mode} --- default lighting, the player-character is solid, movement
  keys move the character.
\item
  \textbf{Insert mode} --- the world desaturates or dims at the edges, the
  character's ``hands'' visibly extend into the floor/wall in front of them
  (they're now writing into the structure itself), and \texttt{hjkl} stop moving the
  camera and start doing nothing (exactly as in real Vim) --- which the game
  should let happen and let feel \emph{wrong}, rather than intercepting it.
\item
  \textbf{Visual mode} (\texttt{v}, \texttt{V}, \texttt{Ctrl-v}) --- a glowing selection volume appears and
  grows as the player moves, previewing exactly what an operator will act on
  before they commit. This is the mode 3D helps most: \texttt{Ctrl-v} (visual block
  mode) is notoriously hard to picture in 2D tutorials and is trivial to
  picture as a literal 3D selection box.
\item
  \textbf{Command mode} (\texttt{:}) --- the world freezes and a console overlay drops down,
  visually distinct from gameplay, reinforcing that \texttt{:} commands are a
  different register of action (save, quit, search-and-replace) from
  buffer editing.
\end{itemize}

\hypertarget{sec-the-verb-noun-grammar-as-the-core-combat-puzzle-loop}{%
\section{The Verb--Noun Grammar as the Core Combat/Puzzle Loop}\label{sec-the-verb-noun-grammar-as-the-core-combat-puzzle-loop}}

Vim's real genius is that operators (\texttt{d}, \texttt{y}, \texttt{c}, \texttt{\textgreater{}}, \texttt{=}) compose with
motions (\texttt{w}, \texttt{\$}, \texttt{\}}, \texttt{f\ x}) and text objects (\texttt{iw}, \texttt{a(}, \texttt{it}) into a small
grammar that generates a huge command space from a few primitives. This
grammar \emph{is} the game's mechanical depth, and should be treated as such rather
than smoothed over:

\begin{itemize}
\tightlist
\item
  \textbf{Operators} are ``verbs'' the player selects or holds (functionally: the
  first keystroke of a pending command). Visually, this could be a glow or
  particle effect anchored to the player, signaling ``an action is armed and
  waiting for a target.''
\item
  \textbf{Motions and text objects} are ``nouns'' --- literally the things in the world
  the glowing verb can be pointed at. \texttt{diw} (delete inner word) becomes: arm
  ``delete,'' point it at a word-object, watch the word-platform vanish and the
  gap close.
\item
  \textbf{Counts} (\texttt{3dw}, \texttt{5j}) become visible charge-up: a numeral appears in the
  HUD as it's typed, and the resulting action visibly happens N times in quick
  succession, so the player \emph{sees} why \texttt{3dw} differs from \texttt{d3w} differs from
  \texttt{3d2w} (they're all the same in real Vim --- 30, and this equivalence is a
  great early puzzle/gotcha for the game to demonstrate rather than explain).
\end{itemize}

This mirrors the same ``verb applied to an addressable target'' grammar already
used elsewhere in your own work --- text objects are exactly a \emph{reach} into a
structured space, the same shape as targeting a coordinate for a rendering
request, just here the space being addressed is a text buffer instead of a
provenance graph.

\hypertarget{sec-advanced-systems}{%
\section{Advanced Systems}\label{sec-advanced-systems}}

\begin{itemize}
\tightlist
\item
  \textbf{Registers} (\texttt{"ayy}, \texttt{"ap}) as an inventory system: named slots the player
  fills by yanking, and can later ``spend'' by pasting from. A level built
  around juggling three registers to rearrange three objects is a natural
  boss puzzle.
\item
  \textbf{Marks} (\texttt{ma}, \texttt{\textasciigrave{}a}) as placeable waypoint beacons the player drops and
  teleports back to --- useful the moment levels get large enough that walking
  back is tedious, which is also exactly when real Vim users start wanting
  marks.
\item
  \textbf{Macros} (\texttt{qa\ ...\ q}, \texttt{@a}, \texttt{10@a}) as a ``record a spell, then recast it''
  mechanic --- record a sequence of actions once, then unleash it N times on a
  row of identical obstacles. This is probably the single most satisfying
  mechanic to render in 3D: watching a recorded macro auto-replay across a
  field of targets is inherently more dramatic in 3D than in a flat grid.
\item
  \textbf{Search} (\texttt{/pattern}, \texttt{n}, \texttt{N}) as a radar/highlighting pulse that lights
  up matching objects across the visible world, with \texttt{n}/\texttt{N} snapping the
  camera to the next/previous hit.
\item
  \textbf{Undo/redo} (\texttt{u}, \texttt{Ctrl-r}) as a visible rewind/fast-forward of the last
  action, not just a state change --- the player should \emph{see} the platform they
  just deleted rebuild itself.
\end{itemize}

\hypertarget{sec-level-world-structure}{%
\section{Level \& World Structure}\label{sec-level-world-structure}}

Following the original's proven shape (new mechanic gated behind a small,
forced-practice arena before it's allowed to mix with prior mechanics), but
organized as a small set of ``documents'' (worlds) rather than one continuous
overworld:

\begin{enumerate}
\def\labelenumi{\arabic{enumi}.}
\tightlist
\item
  \textbf{The Plaza} --- \texttt{hjkl}, \texttt{0}/\texttt{\$}, \texttt{gg}/\texttt{G}. Pure movement, no operators yet.
\item
  \textbf{The Word Foundry} --- \texttt{w}/\texttt{b}/\texttt{e}/\texttt{ge}, introduces the word as the atomic
  object the rest of the grammar will operate on.
\item
  \textbf{The Editing Yards} --- \texttt{i}/\texttt{a}/\texttt{o}/\texttt{O}, first exposure to mode-switching,
  deliberately includes a puzzle solvable only by \emph{escaping} insert mode
  (teaching that \texttt{Esc} is not optional).
\item
  \textbf{The Grammar Works} --- operators (\texttt{d}/\texttt{y}/\texttt{c}) combined with everything
  learned so far. This is where the verb--noun system in Section~\ref{sec-the-verb-noun-grammar-as-the-core-combat-puzzle-loop} becomes
  the primary loop.
\item
  \textbf{Searchworks} --- \texttt{:s}/\texttt{:\%s} and \texttt{:g}/\texttt{:v}, the selection-and-transformation
  family that sits alongside operators and text objects rather than above
  them, since it depends on neither registers, marks, macros, nor undo.
  Placed here rather than later specifically because it doesn't need
  anything the later worlds teach. Every match highlights as a population
  simultaneously, not one at a time, and a substitution or \texttt{:g}-driven
  action resolves that whole population toward its result in one
  coordinated motion --- the world's visual realization of
  match-set → operation → result, letting a player see the population
  before committing to transform it, the same way Visual mode previews an
  operator's effect before it lands.
\item
  \textbf{The Archive Stacks} --- registers, marks, and the buffer-stack Z-axis from
  Option A above; splits and windows become physically real.
\item
  \textbf{The Loop Chamber} --- macros, counts, and repeat (\texttt{.}), building toward
  puzzles only efficiently solvable by recording and replaying a macro.
\item
  \textbf{The Undo Spire} --- the undo-tree-as-tower idea from Option C, offered as
  an optional late-game area for players who want the harder conceptual
  payoff rather than a mandatory main-path level.
\end{enumerate}

\hypertarget{sec-scoring-a-par-system}{%
\section{Scoring: A Par System}\label{sec-scoring-a-par-system}}

Vim Adventures already implicitly rewards fewer keystrokes; making this
explicit as a \textbf{par score per puzzle} (fewest keystrokes to solve, shown
alongside the player's actual count) gives skilled players a reason to replay
early levels with newly learned tricks --- solving Level 1 in 3 keystrokes once
you know \texttt{.} and counts, instead of the 12 it took the first time through, is
its own kind of satisfying. This is the same competitive-efficiency shape as
``pipeline golf'' --- judging a transformation by how few steps it takes to
faithfully reach a target --- just applied to keystrokes instead of rendering
pipelines.

\hypertarget{sec-technical-approach}{%
\section{Technical Approach}\label{sec-technical-approach}}

\begin{itemize}
\tightlist
\item
  \textbf{Engine:} Three.js is the pragmatic choice for a browser-playable
  prototype --- no install friction, and the mechanics here (grid-aligned
  movement, simple selection volumes, particle effects for verbs/macros)
  don't need a heavier engine.
\item
  \textbf{Input layer:} the actual hard part is not the 3D rendering, it's a
  correct modal-keystroke parser --- normal/insert/visual/command mode
  dispatch, pending-operator state, count-prefix accumulation, and the
  motion/text-object grammar. This should be built and tested as a
  standalone, headless module \emph{before} any 3D work starts, so the ``is this
  real Vim behavior'' question is answered independently of how it's rendered.
\item
  \textbf{World representation:} a buffer is a 2D array of characters plus
  metadata (word boundaries, indentation, bracket matching) computed once per
  edit; the 3D scene is a pure function of that array, regenerated on change
  rather than hand-animated per keystroke. This keeps the renderer honest ---
  if the buffer state doesn't change, nothing in the world should either.
\end{itemize}

\hypertarget{sec-open-design-questions}{%
\section{Open Design Questions}\label{sec-open-design-questions}}

\begin{itemize}
\tightlist
\item
  How much of real Vim's \emph{modal escape-hatch chaos} (pressing \texttt{Esc} twice,
  accidentally triggering \texttt{ZZ}, etc.) should the game protect against versus
  let happen? Protecting too much undercuts the ``this is real Vim'' pillar;
  protecting too little risks frustrating brand-new players in a way the
  original 2D game's simpler visual grammar didn't.
\item
  Should text objects beyond \texttt{iw}/\texttt{aw} (e.g.~\texttt{it}/\texttt{at} for tags, \texttt{i"}/\texttt{a"} for
  quotes) require the world to render visible tag/quote structure, or would
  that only make sense in a ``the buffer is actual code'' late-game world
  rather than the more abstract early plaza?
\item
  Multiplayer/co-op is tempting given the macro-replay and register-passing
  mechanics (one player records, another triggers) but adds real complexity
  and isn't needed for the core teaching goal --- worth treating as a stretch
  feature, not a v1 pillar.
\end{itemize}

\hypertarget{sec-what-this-adds-over-the-original}{%
\section{What This Adds Over the Original}\label{sec-what-this-adds-over-the-original}}

Not ``better graphics'' --- the specific things a third dimension makes possible
that a top-down grid cannot: mode as visible world-state rather than a
tracked mental flag, visual-block mode as a literal volume instead of a
description, splits/windows as physically real places instead of a toggled
UI panel, and macros as something dramatic to \emph{watch} replay across a field
of targets rather than a log of keystrokes scrolling past.

\part{Technical Specification}
\setcounter{section}{-1}

\hypertarget{a-3d-spatial-tutorial-for-real-vim-motions-v0.3-merged-with-the-concept-sketch-into-a-single-latex-document}{%
\subsection{A 3D spatial tutorial for real Vim motions --- v0.3, merged with the concept sketch into a single LaTeX document}\label{a-3d-spatial-tutorial-for-real-vim-motions-v0.3-merged-with-the-concept-sketch-into-a-single-latex-document}}

This supersedes the earlier concept sketch's prose sections with precise
specifications, and is merged with it as a single document; heading numbers
below are illustrative in this markdown source and are superseded by
automatic numbering in the merged document.

\begin{center}\rule{0.5\linewidth}{0.5pt}\end{center}

\hypertarget{sec-non-negotiables}{%
\section{Non-Negotiables}\label{sec-non-negotiables}}

These are the correctness bar, not aspirations. A build that violates any of
these has failed the project's actual goal, regardless of how the 3D world
looks.

\begin{enumerate}
\def\labelenumi{\arabic{enumi}.}
\tightlist
\item
  Every keystroke sequence the player can type must produce \textbf{exactly} the
  buffer state real Vim would produce for the same sequence on the same
  input, including the specific off-by-one and edge-case behaviors in
  Section~\ref{sec-motion-operator-semantics-the-part-that-must-be-exact}.
\item
  The input parser (Section~\ref{sec-mode-state-machine}--4) must be implemented and unit-tested as a
  \textbf{headless module with zero rendering dependencies}, before any Three.js
  code is written. The renderer consumes parser output; it never influences
  parser behavior.
\item
  No command, mode, or key that doesn't exist in real Vim is ever bound to
  anything. Cosmetic-only extensions (camera shortcuts, menu keys) must use
  keys Vim itself doesn't bind, and must be documented as out-of-grammar.
\end{enumerate}

\begin{center}\rule{0.5\linewidth}{0.5pt}\end{center}

\hypertarget{sec-formal-command-grammar}{%
\section{Formal Command Grammar}\label{sec-formal-command-grammar}}

A Normal-mode command, in the fragment this project targets for v1, is
generated by the following grammar. This intentionally mirrors the style of
a BNF specification rather than prose, for the same reason a prose
description of Vim's grammar tends to hide exactly the ambiguities that
matter.

\begin{Verbatim}[breaklines=true, breakanywhere=true, fontsize=\small]
<command>       ::= <prefix>* <operator-cmd> | <prefix>* <motion-cmd> | <prefix>* <action>
<prefix>        ::= <count> | "\"" <register>
<operator-cmd>  ::= <operator> <count>? ( <motion> | <text-object> )
                   | <self-apply>                        ; see self-application list below
<motion-cmd>    ::= <motion>
<count>         ::= [1-9][0-9]*
<operator>      ::= "d" | "y" | "c" | ">" | "<" | "=" | "g~" | "gu" | "gU" | "!"
<self-apply>    ::= "dd" | "yy" | "cc" | ">>" | "<<" | "=="
                   | "guu" | "gugu" | "gUU" | "gUgU" | "g~~" | "g~g~"
<motion>        ::= "h" | "l" | "j" | "k"
                   | "0" | "^" | "$" | "gg" | "G"
                   | "w" | "W" | "b" | "B" | "e" | "E" | "ge" | "gE"
                   | "f" <char> | "F" <char> | "t" <char> | "T" <char> | ";" | ","
                   | "{" | "}" | "(" | ")" | "%"
                   | "H" | "M" | "L"
<text-object>   ::= ("i" | "a") <object-kind>
<object-kind>   ::= "w" | "W" | "(" | ")" | "b" | "{" | "}" | "B"
                   | "[" | "]" | "<" | ">" | "\"" | "'" | "t" | "p" | "s"
<action>        ::= "i" | "a" | "I" | "A" | "o" | "O"      ; enter Insert mode
                   | "v" | "V" | "\x16"                    ; enter Visual (char/line/block)
                   | "x" | "X" | "~" | "r" <char>
                   | "u" | "\x12"                          ; undo / redo (Ctrl-r)
                   | "." | "p" | "P"
                   | "q" <register> ... "q"                ; macro record
                   | "@" <register>
                   | "m" <mark> | "`" <mark> | "'" <mark>
                   | "/" <pattern> <CR> | "?" <pattern> <CR> | "n" | "N"
                   | ":" <ex-command> <CR>
\end{Verbatim}

Three properties of this grammar drive most of the parser's complexity and
are worth stating outside the grammar itself:

\begin{itemize}
\tightlist
\item
  \textbf{Register and count prefixes are order-independent.} The previous
  revision of this grammar attached the register prefix only to the
  outside of an already-complete command (\texttt{"a} followed by a whole
  \texttt{\textless{}operator-cmd\textgreater{}}), which cannot generate \texttt{"a3yy} and \texttt{3"ayy} as the same
  command --- real Vim accepts both, with identical effect. The
  \texttt{\textless{}prefix\textgreater{}*} rule above fixes this: \texttt{\textless{}count\textgreater{}} and \texttt{"} \texttt{\textless{}register\textgreater{}} may
  appear in either order, zero or one of each, before the core command.
  If a count appears in a \texttt{\textless{}prefix\textgreater{}} and a second count appears later
  inside \texttt{\textless{}operator-cmd\textgreater{}} (e.g.~\texttt{3"a2yy}), the two counts multiply per the
  existing count-multiplication rule below --- \texttt{3"a2yy} yanks 6 lines into
  register \texttt{a}, the same as \texttt{"a6yy}. Resolution into a \texttt{ResolvedCommand}
  (Section~\ref{sec-command-objects-and-replay-semantics}) discards ordering entirely: the resolved \texttt{count} and
  \texttt{register} fields are the same regardless of which order they were
  typed in.
\item
  \textbf{Count multiplication.} \texttt{\textless{}count\textgreater{}} may appear before the operator \emph{and}
  before the motion in the same command, and the two multiply:
  \texttt{3d2w} deletes \texttt{3\ ×\ 2\ =\ 6} words, identically to \texttt{d6w} or \texttt{6dw}. This is a
  real Vim behavior, not a simplification, and Section~\ref{sec-text-objects} lists it as a
  required test case precisely because it's the first thing an incorrect
  implementation gets wrong.
\item
  \textbf{Operator self-application is linewise, and the v1 grammar fixes the
  full supported set explicitly} rather than leaving it as a ``such as''
  list: \texttt{dd}, \texttt{yy}, \texttt{cc}, \texttt{\textgreater{}\textgreater{}}, \texttt{\textless{}\textless{}}, \texttt{==}, and the case-changing doubles
  \texttt{guu}/\texttt{gugu}, \texttt{gUU}/\texttt{gUgU}, \texttt{g\textasciitilde{}\textasciitilde{}}/\texttt{g\textasciitilde{}g\textasciitilde{}} (both the doubled-letter and
  doubled-operator spellings are accepted for the case operators, matching
  real Vim). Each is grammatically distinct from \texttt{\textless{}operator\textgreater{}\textless{}motion\textgreater{}} and
  always means ``apply linewise to \texttt{\textless{}count\textgreater{}} lines starting at the cursor
  line,'' not ``apply the operator to itself as a motion.'' No other operator
  combination self-applies; \texttt{!!}, \texttt{d=}-style speculative doubles, and
  anything not in \texttt{\textless{}self-apply\textgreater{}} above are simply not valid commands.
\end{itemize}

\begin{center}\rule{0.5\linewidth}{0.5pt}\end{center}

\hypertarget{sec-mode-state-machine}{%
\section{Mode State Machine}\label{sec-mode-state-machine}}

\begin{Verbatim}[breaklines=true, breakanywhere=true, fontsize=\small]
                    ┌─────────────┐
        Esc/Ctrl-[  │             │  i,a,I,A,o,O
       ┌────────────┤   NORMAL    ├────────────┐
       │            │             │            │
       │            └──┬───┬───┬──┘            ▼
       │               │   │   │         ┌───────────┐
       │      v/V/^V   │   │   │  :      │  INSERT   │
       │        ┌──────┘   │   └──────┐  └─────┬─────┘
       │        ▼          │          ▼        │ Esc
       │  ┌───────────┐    │    ┌───────────┐  │
       └──┤  VISUAL   │    │    │  COMMAND  │  │
          │ (3 kinds) │    │    │   LINE    │  │
          └─────┬─────┘    │    └─────┬─────┘  │
                │ operator │          │ CR/Esc │
                │ applied  │          └────────┤
                └──────────┴───────────────────┘
                           │
                    (back to NORMAL)

  Additional transition, not drawn above to keep the box diagram legible:
  VISUAL (any kind) --Esc--> NORMAL, canceling the selection with no
  operator applied and no buffer change. This is a frequent, ordinary
  transition — not an edge case — and is listed here rather than boxed
  above only because a second exit arrow crowded exactly the corner of the
  diagram already carrying the "operator applied" path, making the
  diagram harder to read than it clarified.

  Sub-state of NORMAL, not a separate top-level mode:
  PENDING-OPERATOR — entered after an operator key with no motion yet
  (e.g. "d" alone). Accepts: count, motion, text-object, or the same
  operator again (linewise self-application). Times out to nothing on Esc.
\end{Verbatim}

State is represented as:

\begin{Verbatim}[breaklines=true, breakanywhere=true, fontsize=\small]
type Mode = "NORMAL" | "INSERT" | "VISUAL_CHAR" | "VISUAL_LINE"
          | "VISUAL_BLOCK" | "COMMAND_LINE";

interface ParserState {
  mode: Mode;
  pendingOperator: Operator | null;   // e.g. "d", null if none pending
  pendingCount: number | null;        // count typed before operator
  pendingRegister: string | null;     // "a in "ayy
  countBuffer: string;                // digits typed so far, this command
  visualAnchor: Position | null;      // fixed end of visual selection
  lastFindChar: { cmd: "f"|"F"|"t"|"T"; char: string } | null; // for ; and ,
  lastChange: RecordedCommand | null; // for "."
  recordingMacro: string | null;      // register name, or null
  macroRecording: KeyEvent[];
}
\end{Verbatim}

\begin{center}\rule{0.5\linewidth}{0.5pt}\end{center}

\hypertarget{sec-buffer-model}{%
\section{Buffer Model}\label{sec-buffer-model}}

\begin{Verbatim}[breaklines=true, breakanywhere=true, fontsize=\small]
interface Position { line: number; col: number; }  // 0-indexed internally

interface Buffer {
  lines: string[];                 // one string per line, no trailing \n
  cursor: Position;
  marks: Map<string, Position>;    // 'a'-'z' local, "'" and "`" implicit marks
  undoTree: UndoNode;              // see Section 3.2
  currentUndoNode: UndoNode;
}

interface Register {
  text: string;
  kind: "charwise" | "linewise" | "blockwise";
}
// registers: Map<string, Register>, keys 'a'-'z' (append with 'A'-'Z'),
// '"' (unnamed, implicit target of any yank/delete without explicit register),
// '0' (last yank specifically), '1'-'9' (delete history ring)
\end{Verbatim}

\hypertarget{sec-word-classification}{%
\subsection{Word Classification}\label{sec-word-classification}}

Vim's \texttt{w}/\texttt{b}/\texttt{e} and \texttt{W}/\texttt{B}/\texttt{E} differ in exactly one respect: lowercase
motions treat three character classes as boundaries (\textbf{word-chars}
\texttt{{[}A-Za-z0-9\_{]}}, \textbf{punctuation} \texttt{{[}\^{}A-Za-z0-9\_\textbackslash{}s{]}}, and \textbf{whitespace}), while
uppercase motions collapse word-chars and punctuation into one class
(\textbf{WORD}, anything non-whitespace) and only whitespace is a boundary. This
single classification function is shared by \texttt{w}, \texttt{b}, \texttt{e}, \texttt{ge}, and by the
\texttt{iw}/\texttt{aw} text objects, so it should be implemented exactly once:

\begin{Verbatim}[breaklines=true, breakanywhere=true, fontsize=\small]
type CharClass = "word" | "punct" | "space";
function classify(ch: string): CharClass { /* per rule above */ }
\end{Verbatim}

\hypertarget{sec-undo-tree}{%
\subsection{Undo Tree}\label{sec-undo-tree}}

Vim's undo is a tree, not a stack --- a new edit made after undoing does not
delete the redo branch, it forks it, reachable via \texttt{g-}/\texttt{g+} (time-travel
through \emph{all} states in the order visited) or \texttt{u}/\texttt{Ctrl-r} (parent/last-child
only). This is exactly the ``World D --- Undo Spire'' content, and the data
structure should be built as a tree from the start rather than retrofitted:

\begin{Verbatim}[breaklines=true, breakanywhere=true, fontsize=\small]
interface UndoNode {
  id: number;
  bufferSnapshot: string[];   // full snapshot; diffing is an optimization, not v1
  parent: UndoNode | null;
  children: UndoNode[];
  seq: number;                // global sequence number, for g-/g+ ordering
  timestamp: number;
}
\end{Verbatim}

\hypertarget{sec-command-objects-and-replay-semantics}{%
\subsection{Command Objects and Replay Semantics}\label{sec-command-objects-and-replay-semantics}}

Macros and \texttt{.} look similar from the outside --- both ``do something again'' --- but
they replay at two different levels of representation, and conflating them
is a reliable source of bugs.

\textbf{\texttt{.} replays a resolved command, not keystrokes.} The parser must retain a
structured record of the last change, not the raw keys that produced it.
The previous revision of this document gave \texttt{ResolvedCommand} a single
shape built around \texttt{operator} and \texttt{target}, which has no room for
count-prefixed insertions (\texttt{3ohello\textless{}Esc\textgreater{}}, \texttt{3ihello\textless{}Esc\textgreater{}}) or simple
repeated actions (\texttt{3x}, \texttt{5\textasciitilde{}}) --- neither has an operator or a
motion/text-object target in the sense Section~\ref{sec-formal-command-grammar}'s grammar defines. The
corrected shape is a tagged union:

\begin{Verbatim}[breaklines=true, breakanywhere=true, fontsize=\small]
type ResolvedCommand =
  | {
      kind: "operator-motion";
      operator: Operator;
      count: number;              // fully multiplied — 6, not "3" and "2" separately
      target: Motion | TextObject;
      insertedText?: string;      // present only for c-style changes
      register?: string;
    }
  | {
      kind: "insert-action";
      action: "i" | "a" | "I" | "A" | "o" | "O";
      count: number;
      insertedText: string;       // captured verbatim between entering Insert and Esc
    }
  | {
      kind: "simple-action";
      action: "x" | "X" | "~" | "r" | "J" | "gJ";
      count: number;
      char?: string;               // present only for "r"
      register?: string;
    };
\end{Verbatim}

\texttt{3d2w} and \texttt{d6w} resolve to the identical \texttt{\{kind:\ "operator-motion",\ operator:\ "d",\ count:\ 6,\ target:\ "w"\}}, and it is this resolved form --- not
either surface spelling --- that \texttt{.} replays, at whatever position the
cursor currently occupies. This is also why \texttt{.} after a \texttt{c}-style change
retypes \texttt{insertedText} verbatim rather than re-entering Insert mode for the
player to retype.

\textbf{Count semantics differ by insert-action, and this difference is not
optional to get right.} \texttt{i}, \texttt{a}, \texttt{I}, and \texttt{A} repeat the \emph{inserted text in
place}: \texttt{3ihello\textless{}Esc\textgreater{}} on an empty line produces a single line reading
\texttt{hellohellohello} --- the count multiplies how many times the same text is
appended at the insertion point, not how many separate edits occur. \texttt{o} and
\texttt{O} repeat the \emph{whole open-line-and-insert operation}: \texttt{3ohello\textless{}Esc\textgreater{}}
produces three new lines, each independently containing \texttt{hello}, because
each repetition opens its own new line before inserting. A \texttt{ResolvedCommand}
of \texttt{kind:\ "insert-action"} must therefore branch on \texttt{action} at replay
time --- the same \texttt{count} field means two structurally different replays
depending on whether \texttt{action} is in \texttt{\{i,\ a,\ I,\ A\}} or \texttt{\{o,\ O\}} --- rather than
being interpreted uniformly the way \texttt{count} is for \texttt{operator-motion}.

\textbf{\texttt{J}'s count means something different from every other simple action's
count.} For \texttt{x}, \texttt{X}, \texttt{\textasciitilde{}}, and \texttt{r}, \texttt{count} means ``repeat this action N
times'' in the ordinary sense (\texttt{3x} deletes 3 characters). For \texttt{J}, \texttt{count}
means ``join this many lines total, counting the current line'' --- \texttt{3J} joins
the current line with the next 2, not the current line with the next 3.
This mirrors \texttt{G}'s count meaning ``go to line N'' rather than ``repeat G N
times'' (Section~\ref{sec-motion-operator-semantics-the-part-that-must-be-exact}), and should be implemented as the same kind of named
special case, not derived from the general repeat-N-times rule.

\textbf{A new count given at \texttt{.}-time replaces, not multiplies.} If the player
types a count immediately before \texttt{.} (e.g.~\texttt{5.} following an earlier \texttt{3dw}),
the stored count is \emph{replaced} by 5 for this one replay --- the result is
\texttt{d5w}, not \texttt{d15w}, and the replacement does not persist: a bare \texttt{.}
afterward reuses whichever count last stood, following the same
replace-not-accumulate rule. This is a real, easily-missed Vim behavior and
belongs in the Section~\ref{sec-testing-plan} test suite as its own case.

\textbf{Macros replay raw keystrokes, not resolved commands.} A macro register
(\texttt{qa\ ...\ q}) stores the literal \texttt{KeyEvent} sequence typed during recording,
and \texttt{@a} re-feeds that sequence through the ordinary parser from scratch ---
it does not store or replay \texttt{ResolvedCommand} objects. This is why a macro
built around \texttt{dw} deletes whatever word is actually under the cursor at each
replay site, rather than a fixed span: the macro has no memory of what the
first replay resolved to, only of which keys were pressed. \texttt{10@a} repeats
the whole recorded sequence ten times, not any single command within it six
or ten times.

\hypertarget{sec-cursor-invariants}{%
\subsection{Cursor Invariants}\label{sec-cursor-invariants}}

Most bugs in a from-scratch Vim clone are cursor bugs, not edit bugs --- the
edit itself lands correctly and the cursor ends up somewhere Vim would never
put it. The invariants below should be enforced as a single clamping
function called after every buffer mutation, not re-derived ad hoc per
command.

\begin{enumerate}
\def\labelenumi{\arabic{enumi}.}
\tightlist
\item
  \texttt{cursor.line} is always in \texttt{{[}0,\ lines.length\ -\ 1{]}}.
\item
  In \texttt{NORMAL}, \texttt{VISUAL\_*}, and \texttt{COMMAND\_LINE} modes, \texttt{cursor.col} is
  clamped to \texttt{{[}0,\ max(0,\ lineLength\ -\ 1){]}} --- the cursor may never rest one
  past the last character of a non-empty line. An empty line clamps to
  column 0.
\item
  In \texttt{INSERT} mode only, \texttt{cursor.col} may equal \texttt{lineLength} (one position
  past the last character), since that is where typed text would be
  appended. Leaving Insert mode re-applies rule 2 immediately, which is why
  the cursor visibly steps back one column on \texttt{Esc} at end-of-line --- this
  is correct behavior to replicate, not a bug to fix.
\item
  After a linewise delete, the cursor is not preserved at its prior column;
  it moves to the first non-blank column of whichever line now occupies
  the cursor's old line index.
\item
  After a charwise delete that removes the character the cursor was on,
  the cursor clamps left under rule 2 if the deletion shortened the line
  below the old column.
\item
  \texttt{J} (join) leaves the cursor at the join point --- the single space
  inserted between the two joined lines (or, for \texttt{gJ}, at the exact
  boundary, since \texttt{gJ} does not insert a space).
\item
  A mark referencing a line that has since been deleted is cleared, not
  left dangling. Jumping to a cleared mark (\texttt{\textasciigrave{}a} or \texttt{\textquotesingle{}a}) is a
  documented no-op, not an error condition the parser needs to surface.
\end{enumerate}

\begin{center}\rule{0.5\linewidth}{0.5pt}\end{center}

\hypertarget{sec-motion-operator-semantics-the-part-that-must-be-exact}{%
\section{Motion \& Operator Semantics (the part that must be exact)}\label{sec-motion-operator-semantics-the-part-that-must-be-exact}}

Every motion has a \textbf{type} that determines how an operator combined with it
behaves. Getting this table wrong is the single most common way a
from-scratch Vim clone feels almost-but-not-quite right.

\begin{Verbatim}[breaklines=true, breakanywhere=true, fontsize=\small]
type MotionKind = "exclusive" | "inclusive" | "linewise";
\end{Verbatim}

\begin{itemize}
\tightlist
\item
  \textbf{Exclusive} --- the operator's range is \texttt{{[}start,\ destination)}: the
  character the motion lands the cursor on is not included.
\item
  \textbf{Inclusive} --- the range is \texttt{{[}start,\ destination{]}}: the destination
  character is included.
\item
  \textbf{Linewise} --- the range expands to whole lines regardless of the columns
  either endpoint's motion actually computed, from the first line touched
  through the last, newline characters included on both ends.
\end{itemize}

\begingroup\small
\begin{longtable}[]{@{}
  >{\raggedright\arraybackslash}p{(\columnwidth - 4\tabcolsep) * \real{0.3333}}
  >{\raggedright\arraybackslash}p{(\columnwidth - 4\tabcolsep) * \real{0.3333}}
  >{\raggedright\arraybackslash}p{(\columnwidth - 4\tabcolsep) * \real{0.3333}}@{}}
\toprule\noalign{}
\begin{minipage}[b]{\linewidth}\raggedright
Motion
\end{minipage} & \begin{minipage}[b]{\linewidth}\raggedright
Type
\end{minipage} & \begin{minipage}[b]{\linewidth}\raggedright
Notes
\end{minipage} \\
\midrule\noalign{}
\endhead
\bottomrule\noalign{}
\endlastfoot
\texttt{h}, \texttt{l} & exclusive & won't cross line boundaries \\
\texttt{j}, \texttt{k} & linewise & \\
\texttt{0} & exclusive & \\
\texttt{\^{}} & exclusive & first non-blank \\
\texttt{\$} & inclusive & \textbf{special-cased}: \texttt{d\$} deletes to end of line \emph{and} the type is inclusive, unlike most exclusive-by-default column motions \\
\texttt{w} & exclusive & \textbf{special-cased}: if \texttt{w} is the object of an operator and lands on the first non-blank of a new line, it's pulled back to end of previous line (\texttt{dw} at end of line does not delete the newline) \\
\texttt{W} & exclusive & same special case as \texttt{w} \\
\texttt{b}, \texttt{B} & exclusive & \\
\texttt{e}, \texttt{E} & inclusive & \\
\texttt{ge}, \texttt{gE} & inclusive & \\
\texttt{f\{c\}}, \texttt{t\{c\}} & inclusive & \\
\texttt{F\{c\}}, \texttt{T\{c\}} & exclusive & \\
\texttt{gg}, \texttt{G}, \texttt{\{count\}G} & linewise & \\
\texttt{\{}, \texttt{\}} & exclusive & \\
\texttt{(}, \texttt{)} & exclusive & \\
\texttt{\%} & inclusive & \\
\texttt{H}, \texttt{M}, \texttt{L} & linewise & \\
\end{longtable}
\endgroup

\textbf{The \texttt{cw}/\texttt{ce} special case.} By Vim's documented behavior, if the cursor
is on a non-blank character, \texttt{cw} behaves like \texttt{ce} (changes to the end of
the current word, not to the start of the next one) --- the one deliberate
exception carved out of the ``operator + motion'' composition rule, added
specifically because \texttt{dw}'s ``don't eat the trailing whitespace into the next
word'' behavior felt wrong for \texttt{c}. This must be implemented as an explicit
special case, not derived from the general rule, because it isn't derivable
from it.

\textbf{Linewise operator range.} When an operator combines with a linewise
motion (\texttt{dj}, \texttt{d\}} is \emph{not} linewise --- \texttt{\}} is exclusive; \texttt{dG}, \texttt{dgg} \emph{are}),
the operator acts on whole lines, including the newline character, and the
cursor lands on the first non-blank of the resulting line --- not at the
column it started on.

\textbf{Required test cases} (Section~\ref{sec-testing-plan} formalizes this as an actual test
suite; this is the minimum set that must pass before any 3D work begins):

\begin{Verbatim}[breaklines=true, breakanywhere=true, fontsize=\small]
buffer: "The quick brown fox"     cursor col 0
  dw   → "quick brown fox"
  de   → "The quick brown fox" minus "The" but with trailing space (differs from dw!)
  cw   → behaves as ce: enters insert with " quick brown fox" and cursor after "The"

buffer: "foo"                      cursor col 0  (word at end of buffer, no trailing word)
  dw   → "" (does not error, does not hang)

buffer: "foo\nbar"                 cursor at end of "foo"
  dw   → deletes to end of "foo" only; does NOT pull "bar" up (newline preserved)

count multiplication:
  3d2w  ==  d6w  ==  6dw   (must all produce identical results)

linewise self-application:
  3dd   deletes 3 lines starting at cursor line, cursor lands at first non-blank
        of the line that shifts into the cursor's old position
\end{Verbatim}

\begin{center}\rule{0.5\linewidth}{0.5pt}\end{center}

\hypertarget{sec-text-objects}{%
\section{Text Objects}\label{sec-text-objects}}

Text objects only ever appear as the target of an operator or in Visual
mode; they are never standalone motions. \texttt{i} (``inner'') excludes the
delimiter/surrounding whitespace; \texttt{a} (``a'', i.e.~``around'') includes it.

\begingroup\small
\begin{longtable}[]{@{}
  >{\raggedright\arraybackslash}p{(\columnwidth - 4\tabcolsep) * \real{0.3333}}
  >{\raggedright\arraybackslash}p{(\columnwidth - 4\tabcolsep) * \real{0.3333}}
  >{\raggedright\arraybackslash}p{(\columnwidth - 4\tabcolsep) * \real{0.3333}}@{}}
\toprule\noalign{}
\begin{minipage}[b]{\linewidth}\raggedright
Object
\end{minipage} & \begin{minipage}[b]{\linewidth}\raggedright
\texttt{i} behavior
\end{minipage} & \begin{minipage}[b]{\linewidth}\raggedright
\texttt{a} behavior
\end{minipage} \\
\midrule\noalign{}
\endhead
\bottomrule\noalign{}
\endlastfoot
\texttt{iw} / \texttt{aw} & current word, no surrounding space & word plus one side of adjacent whitespace \\
\texttt{i(} \texttt{i)} \texttt{ib} & contents between nearest enclosing \texttt{(} \texttt{)} & includes the parens themselves \\
\texttt{i\{} \texttt{i\}} \texttt{iB} & contents between nearest enclosing \texttt{\{} \texttt{\}} & includes the braces \\
\texttt{i"} / \texttt{a"} & contents between nearest quote pair on the line & includes the quotes \\
\texttt{it} / \texttt{at} & contents between nearest enclosing HTML/XML tag pair & includes the tags \\
\texttt{ip} / \texttt{ap} & current paragraph (blank-line delimited) & paragraph plus one trailing blank line \\
\end{longtable}
\endgroup

Nesting matters for \texttt{i(}/\texttt{a(} etc.: with the cursor inside nested
parentheses, the object is always the \emph{nearest enclosing} pair, and
\texttt{2i(} (a count before a text object) walks outward one level per count.

\hypertarget{sec-text-object-resolution}{%
\subsection{Text Object Resolution}\label{sec-text-object-resolution}}

Text object resolution is not one uniform algorithm --- it is two, selected
by object kind, and the spec should say so explicitly rather than imply a
single flowchart covers every row of the table above.

\textbf{Bracket and tag objects} (\texttt{(} \texttt{)} \texttt{b}, \texttt{\{} \texttt{\}} \texttt{B}, \texttt{{[}} \texttt{{]}}, \texttt{\textless{}} \texttt{\textgreater{}}, \texttt{t})
resolve by nested balanced search:

\begin{enumerate}
\def\labelenumi{\arabic{enumi}.}
\tightlist
\item
  If the cursor sits on the opening or closing delimiter itself, that pair
  is treated as the innermost enclosing pair.
\item
  Otherwise, search outward from the cursor for the nearest pair of
  matching delimiters that encloses the cursor's position; this search may
  cross line boundaries.
\item
  Apply the count: each additional count expands one level further
  outward (\texttt{2i(} means the pair enclosing the pair found in step 2).
\item
  Apply \texttt{i}/\texttt{a}: \texttt{a} includes the delimiters themselves; \texttt{i} excludes them
  and, in the documented special case where the interior spans multiple
  full lines on its own, also trims the newline immediately inside each
  delimiter, yielding a linewise rather than charwise range for that case
  specifically.
\end{enumerate}

\textbf{Quote objects} (\texttt{i"} \texttt{a"}, \texttt{i\textquotesingle{}} \texttt{a\textquotesingle{}}, \texttt{i\textasciigrave{}} \texttt{a\textasciigrave{}}) resolve
differently, and this difference is not an oversight to reconcile with the
bracket algorithm:

\begin{enumerate}
\def\labelenumi{\arabic{enumi}.}
\tightlist
\item
  The search is restricted to the cursor's current line only --- quotes are
  not treated as nesting across lines the way brackets do.
\item
  If the cursor is already inside a quoted span on that line, that span is
  the object.
\item
  If the cursor sits before the first quote on the line, Vim selects the
  \emph{first} quoted string on the line, rather than searching backward from
  the cursor.
\item
  There is no expand-outward behavior for counts on quote objects, since
  there is nothing for a count to expand into --- a line has no nesting of
  quoted spans in the sense brackets nest.
\end{enumerate}

\begin{center}\rule{0.5\linewidth}{0.5pt}\end{center}

\hypertarget{sec-rendering-mapping-renderer-consumes-parser-state-never-the-reverse}{%
\section{Rendering Mapping (renderer consumes parser state; never the reverse)}\label{sec-rendering-mapping-renderer-consumes-parser-state-never-the-reverse}}

\hypertarget{sec-fidelity-principle}{%
\subsection{Fidelity Principle}\label{sec-fidelity-principle}}

The world exists solely to visualize parser state. Where visual intuition
and actual Vim behavior conflict, Vim behavior wins, unconditionally --- a
world redesign to make some interaction feel more natural is never a
license to change what a keystroke does. Data flows one way:

\begin{Verbatim}[breaklines=true, breakanywhere=true, fontsize=\small]
Parser  →  BufferModel / ParserState  →  Renderer
\end{Verbatim}

and never the reverse. No rendering-layer code may write to \texttt{ParserState},
\texttt{Buffer}, or \texttt{UndoNode}; the renderer's only inputs are read-only snapshots
of parser output. This restates Section~\ref{sec-non-negotiables}'s second non-negotiable as a
concrete data-flow rule specifically so that a violation is visible on
sight in code review --- a pull request containing a call from renderer code
into parser or buffer state is wrong by inspection, without needing the
Section~\ref{sec-testing-plan} test suite to catch it at runtime.

\begingroup\small
\begin{longtable}[]{@{}
  >{\raggedright\arraybackslash}p{(\columnwidth - 2\tabcolsep) * \real{0.5000}}
  >{\raggedright\arraybackslash}p{(\columnwidth - 2\tabcolsep) * \real{0.5000}}@{}}
\toprule\noalign{}
\begin{minipage}[b]{\linewidth}\raggedright
Parser state
\end{minipage} & \begin{minipage}[b]{\linewidth}\raggedright
Visual representation
\end{minipage} \\
\midrule\noalign{}
\endhead
\bottomrule\noalign{}
\endlastfoot
\texttt{mode:\ NORMAL} & full lighting, player mesh solid and idle-animated \\
\texttt{mode:\ INSERT} & scene desaturates \textasciitilde30\%, player mesh extends a ``writing'' appendage into the nearest surface, \texttt{hjkl} visibly do nothing when pressed (no camera response) \\
\texttt{mode:\ VISUAL\_CHAR/LINE/BLOCK} & translucent selection volume grows live from \texttt{visualAnchor} to cursor; block mode renders as an actual rectangular prism, not a flat highlight \\
\texttt{pendingOperator\ ≠\ null} & particle glow around player, color-coded per operator (\texttt{d} red, \texttt{y} blue, \texttt{c} yellow), persists until a motion/text-object resolves it or Esc cancels it \\
\texttt{countBuffer\ ≠\ ""} & numeral rendered above player HUD-style, not in-world \\
\texttt{recordingMacro\ ≠\ null} & a visibly rotating recording icon above player; on \texttt{@a} playback, the recorded action sequence visibly re-executes at accelerated but perceptible speed across each target in sequence \\
mark set (\texttt{ma}) & a small labeled beacon placed at that world position, persistent until overwritten \\
search active (\texttt{/pattern}) & all matching objects pulse/outline in the world simultaneously; \texttt{n}/\texttt{N} snaps camera focus to next/previous match in document order \\
undo/redo & the affected region visibly reverts/reapplies over \textasciitilde200ms, not an instant cut, so the causal link between \texttt{u} and the specific change is legible \\
\end{longtable}
\endgroup

\begin{center}\rule{0.5\linewidth}{0.5pt}\end{center}

\hypertarget{sec-level-data-format}{%
\section{Level Data Format}\label{sec-level-data-format}}

\begin{Verbatim}[breaklines=true, breakanywhere=true, fontsize=\small]
{
  "id": "grammar-works-03",
  "title": "The Grammar Works — Room 3",
  "initialBuffer": ["The quick brown fox", "jumps over the lazy dog"],
  "initialCursor": { "line": 0, "col": 0 },
  "allowedCommands": ["h","l","j","k","w","b","e","d","y","c","p","P","0","$"],
  "goal": {
    "type": "bufferEquals",
    "target": ["quick brown", "jumps over the lazy dog"]
  },
  "par": 5,
  "hints": [
    { "afterKeystrokes": 15, "text": "You don't need to delete word by word." }
  ]
}
\end{Verbatim}

\texttt{allowedCommands} gates which keys the parser will accept in this level
(unrecognized-but-otherwise-valid Vim commands should give an in-world
``not yet unlocked'' cue, not silently fail) --- this is how the progressive
teaching curve from the concept doc is actually enforced, rather than left
as a level-design convention.

\texttt{goal.type} should support at minimum \texttt{bufferEquals} (exact match),
\texttt{bufferMatches} (regex, for goals with multiple valid solutions), and
\texttt{cursorAt} (motion-only levels where the buffer never changes).

\hypertarget{sec-pedagogical-layer}{%
\subsection{Pedagogical Layer}\label{sec-pedagogical-layer}}

\texttt{allowedCommands} implies a binary available/unavailable split; the actual
teaching design needs a third state, and the three should not be
implemented the same way:

\begin{Verbatim}[breaklines=true, breakanywhere=true, fontsize=\small]
type CommandAvailability = "unsupported" | "locked" | "unlocked";
\end{Verbatim}

\begin{itemize}
\tightlist
\item
  \textbf{\texttt{unsupported}} --- not implemented by the parser in this version at all
  (see Section~\ref{sec-future-vim-coverage-explicitly-out-of-scope-for-v1}'s roadmap). Typing it is a no-op with a generic
  not-available cue.
\item
  \textbf{\texttt{locked}} --- fully implemented and semantically correct in the parser,
  but withheld from the player by the current level's \texttt{allowedCommands}.
  Critically, a locked command must still be \emph{parsed and resolved in full}
  by the engine, and only then have its resulting state change discarded in
  favor of a specific educational cue (``you'll unlock delete-word soon'') ---
  never short-circuited before parsing. Implementing it this way keeps the
  parser's grammar coverage complete and independently testable against
  Section~\ref{sec-testing-plan}'s oracle even for commands no level has unlocked yet, rather
  than leaving locked commands as untested dead code until the level that
  first exposes them.
\end{itemize}

\textbf{Gating happens at the resolved-command level, not at the keystroke
level, and this is a decision rather than an implementation detail.}
Individual key-presses are never blocked or intercepted as they're typed ---
a player can freely type any sequence Section~\ref{sec-formal-command-grammar}'s grammar accepts,
including the leading keys of a locked command, exactly as they could in
real Vim. The check happens once a full \texttt{ResolvedCommand} (Section~\ref{sec-command-objects-and-replay-semantics})
exists: only then is its availability looked up, and only a \texttt{locked} or
\texttt{unsupported} result triggers the discard-and-cue behavior described above.
The alternative --- refusing individual keystrokes that could only ever lead
to a locked command --- was considered and rejected, because it would mean
the game's input handling diverges from Section~\ref{sec-formal-command-grammar}'s grammar depending on
which level is loaded, which is exactly the kind of renderer/level-specific
special-casing Section~\ref{sec-fidelity-principle}'s Fidelity Principle rules out for rendering and
should equally be ruled out for input handling.
- \textbf{\texttt{unlocked}} --- available, and behaves exactly per Sections~\ref{sec-formal-command-grammar}--\ref{sec-text-objects}, with no
further distinction from ordinary Vim behavior.

\begin{center}\rule{0.5\linewidth}{0.5pt}\end{center}

\hypertarget{sec-testing-plan}{%
\section{Testing Plan}\label{sec-testing-plan}}

The parser module ships with a table-driven test suite before any rendering
work starts:

\begin{Verbatim}[breaklines=true, breakanywhere=true, fontsize=\small]
interface ParserTestCase {
  initialBuffer: string[];
  initialCursor: Position;
  keys: string;              // raw keystrokes, e.g. "3dw"
  expectedBuffer: string[];
  expectedCursor: Position;
  expectedMode?: Mode;       // default NORMAL
}
\end{Verbatim}

The Section~\ref{sec-motion-operator-semantics-the-part-that-must-be-exact} ``required test cases'' become the seed set; coverage should
expand to include every row of the Section~\ref{sec-motion-operator-semantics-the-part-that-must-be-exact} and Section~\ref{sec-text-objects} tables at least
once, plus the count-multiplication identities as property-based tests
(generate random \texttt{(a,\ b)} and assert \texttt{a*b\ da*b\ w} ≡ \texttt{d(a*b)w} for small
values) rather than a handful of hand-picked examples.

\hypertarget{sec-differential-testing-against-real-vim}{%
\subsection{Differential Testing Against Real Vim}\label{sec-differential-testing-against-real-vim}}

Hand-written expected outputs, however careful, will not catch the edge
cases that matter most, precisely because the person writing the test case
and the person writing the parser share the same blind spots. The stronger
design is an oracle: run the same buffer and keystrokes through real Vim
semantics and assert Buffer World's parser agrees.

The concrete mechanism should be \textbf{Neovim's \texttt{-\/-embed} RPC API}, not a
shelled-out terminal Vim with output scraped from a pseudo-terminal. The
RPC approach avoids ANSI/terminal-emulation fragility entirely and gives
exact, structured buffer and cursor state directly:

\begin{Verbatim}[breaklines=true, breakanywhere=true, fontsize=\small]
interface OracleTestCase {
  initialBuffer: string[];
  initialCursor: Position;
  keys: string;   // raw keystrokes, exactly as Section 1's grammar would parse them
}

async function runOracle(tc: OracleTestCase): Promise<{ buffer: string[]; cursor: Position }> {
  // 1. spawn `nvim --embed --headless`
  // 2. attach via the msgpack-RPC client (the `neovim` npm package)
  // 3. nvim_buf_set_lines(0, 0, -1, false, tc.initialBuffer)
  // 4. nvim_win_set_cursor(0, [tc.initialCursor.line + 1, tc.initialCursor.col])  // nvim rows are 1-indexed
  // 5. nvim_input(tc.keys)
  // 6. read back nvim_buf_get_lines + nvim_win_get_cursor
  // 7. shut the embedded instance down
}
\end{Verbatim}

Given this harness, the test suite becomes generative rather than purely
curated: run a large corpus of keystroke sequences --- both the curated cases
from Section~\ref{sec-motion-operator-semantics-the-part-that-must-be-exact}/5 and randomly generated sequences over a range of seed
buffers --- through both \texttt{runOracle} and Buffer World's own parser, and assert
equality. Any mismatch is either a genuine parser bug, or a deliberate,
documented divergence where v1 intentionally doesn't implement something
real Vim does (an \texttt{unsupported} command per Section~\ref{sec-pedagogical-layer} and Section~\ref{sec-future-vim-coverage-explicitly-out-of-scope-for-v1}). The
harness needs a way to mark specific oracle mismatches as known and
intentional rather than treating 100\% oracle agreement as the bar for every
possible input, since that bar is equivalent to reimplementing all of Vim.

This is a \textbf{Phase 0 requirement}, not a later hardening pass: Section~\ref{sec-build-phases}'s
gate from Phase 0 to Phase 1 should require the oracle suite passing (net of
documented divergences) alongside the hand-written Section~\ref{sec-testing-plan} cases, since
oracle testing is what catches the mismatches a spec document --- including
this one --- cannot anticipate by construction.

\begin{center}\rule{0.5\linewidth}{0.5pt}\end{center}

\hypertarget{sec-build-phases}{%
\section{Build Phases}\label{sec-build-phases}}

The ordering below follows increasing abstraction rather than the order
features were specified in this document: move, edit, select, transform,
store, repeat, time-travel, create. Selection-and-transformation
(\texttt{:s}/\texttt{:g}, Phase 4) is placed immediately after Grammar Works and before
registers, marks, macros, or the undo tree, because it does not
conceptually depend on any of them --- it is a second grammar family
alongside operators and text objects, not an advanced feature layered on
top of the systems that follow it.

\begingroup\small
\begin{longtable}[]{@{}
  >{\raggedright\arraybackslash}p{(\columnwidth - 4\tabcolsep) * \real{0.3333}}
  >{\raggedright\arraybackslash}p{(\columnwidth - 4\tabcolsep) * \real{0.3333}}
  >{\raggedright\arraybackslash}p{(\columnwidth - 4\tabcolsep) * \real{0.3333}}@{}}
\toprule\noalign{}
\begin{minipage}[b]{\linewidth}\raggedright
Phase
\end{minipage} & \begin{minipage}[b]{\linewidth}\raggedright
Deliverable
\end{minipage} & \begin{minipage}[b]{\linewidth}\raggedright
Gate to next phase
\end{minipage} \\
\midrule\noalign{}
\endhead
\bottomrule\noalign{}
\endlastfoot
0 & Headless parser + full Section~\ref{sec-testing-plan} test suite passing, including the Section~\ref{sec-differential-testing-against-real-vim} oracle harness & 100\% of Section~\ref{sec-motion-operator-semantics-the-part-that-must-be-exact}/5 required cases pass, and the oracle suite passes net of documented divergences \\
1 & Three.js renderer, Plaza world only (\texttt{hjkl}, \texttt{0}, \texttt{\$}, \texttt{gg}, \texttt{G}) & movement feels correct, no operators yet \\
2 & Word Foundry + Editing Yards (\texttt{w/b/e}, \texttt{i/a/o/O}, Esc) & mode visualization (Section~\ref{sec-rendering-mapping-renderer-consumes-parser-state-never-the-reverse}) implemented \\
3 & Grammar Works (\texttt{d/y/c} × all Section~\ref{sec-motion-operator-semantics-the-part-that-must-be-exact} motions + Section~\ref{sec-text-objects} text objects) & full operator-motion grammar passes in-world \\
4 & Searchworks (\texttt{:s}, \texttt{:\%s}, \texttt{:g}/\texttt{:v} per Section~\ref{sec-global-commands-and-substitute}) & Section~\ref{sec-global-commands-and-substitute}'s grammar passes its own test cases and oracle agreement, independent of registers/macros/undo \\
5 & Archive Stacks (registers, marks, \texttt{Ctrl-w} splits as Z-axis) & Option A from the concept doc realized \\
6 & Loop Chamber (macros, counts, \texttt{.} repeat) & macro replay visualization (Section~\ref{sec-rendering-mapping-renderer-consumes-parser-state-never-the-reverse}) shipped \\
7 & Undo Spire (Section~\ref{sec-undo-tree} tree, \texttt{g-}/\texttt{g+}) & Section~\ref{sec-undo-tree}'s tree structure and \texttt{g-}/\texttt{g+} ordering verified against the oracle \\
8 & User-generated/editable levels (Section~\ref{sec-user-generated-and-editable-levels}), using the Section~\ref{sec-global-commands-and-substitute} global/substitute grammar & leaderboard verification (13.3) passes against the Section~\ref{sec-differential-testing-against-real-vim} oracle harness \\
\end{longtable}
\endgroup

\hypertarget{sec-searchworks}{%
\subsection{Searchworks}\label{sec-searchworks}}

Searchworks is the first world built around a different mental model than
everything before it. Plaza through Grammar Works are all \emph{move-and-act-on-
one-thing-at-a-time}; Searchworks introduces \emph{select-a-population, then
transform it}. The visual language should mark this shift rather than
reuse the single-target glow from Section~\ref{sec-rendering-mapping-renderer-consumes-parser-state-never-the-reverse}'s operator visualization:

\begin{itemize}
\tightlist
\item
  A search or pattern (\texttt{/pattern}, or the search side of \texttt{:s}/\texttt{:g})
  highlights every matching object across the visible world
  \textbf{simultaneously}, not one at a time --- this is the population being
  selected, and it should read as a population, not a sequence of
  individually-found targets.
\item
  A substitution or \texttt{:g}-driven action then visibly resolves that entire
  highlighted population toward a fixed point in one coordinated motion ---
  every matching object transforming together --- rather than replaying the
  same animation once per match in sequence, which would misrepresent
  \texttt{:g}'s actual scan-then-act-once semantics from Section~\ref{sec-the-global-command-family} as something
  closer to a loop the player is watching iterate.
\item
  This is the concrete world-level realization of the
  match-set → operation → result model that Section~\ref{sec-global-commands-and-substitute} specifies formally:
  the world shows the match set as a visible population before any operation
  is applied, exactly so that the player can predict the result before
  committing to \texttt{:s} or \texttt{:g}, the same way Visual mode's selection volume
  (Section~\ref{sec-rendering-mapping-renderer-consumes-parser-state-never-the-reverse}) lets a player preview an operator's effect before applying it.
\end{itemize}

Par scores and per-level best-keystroke tracking (Section~\ref{sec-level-data-format}'s \texttt{par} field)
can be persisted via the artifact storage API discussed for other projects
(\texttt{window.storage}, keyed per level id) once this moves from a standalone
prototype into a Claude-artifact build --- worth deferring until Phase 1 is
playable, since a persistence layer with no game to persist state for is
premature.

\begin{center}\rule{0.5\linewidth}{0.5pt}\end{center}

\hypertarget{sec-future-vim-coverage-explicitly-out-of-scope-for-v1}{%
\section{Future Vim Coverage (Explicitly Out of Scope for v1)}\label{sec-future-vim-coverage-explicitly-out-of-scope-for-v1}}

This section exists so that a missing feature reads as a decision, not an
oversight. Anything not named in Sections~\ref{sec-formal-command-grammar}--\ref{sec-testing-plan} or listed below as a future
family is out of scope for v1, and a reviewer or contributor proposing it
should check this list before assuming it was simply forgotten.

\begin{itemize}
\tightlist
\item
  \textbf{Window and tab management beyond Section~\ref{sec-buffer-model}'s Z-axis splits} --- \texttt{:tabnew},
  tab navigation, and window-layout commands beyond the basic split model
  Phase 5 implements.
\item
  \textbf{Replace mode (\texttt{R})} --- sticky multi-character replace, where typed
  characters overwrite existing ones until \texttt{Esc}, is excluded for v1. This
  is a deliberate exclusion, not an omission: Section~\ref{sec-mode-state-machine}'s \texttt{Mode} type has no
  \texttt{REPLACE} state, and adding one would touch the state machine, the cursor
  invariants (Section~\ref{sec-cursor-invariants}), and the rendering mapping (Section~\ref{sec-rendering-mapping-renderer-consumes-parser-state-never-the-reverse}) for a
  single mode family that isn't central to the curriculum this document
  specifies. This is distinct from single-character replace, \texttt{r\textless{}char\textgreater{}},
  which is already in \texttt{\textless{}action\textgreater{}} (Section~\ref{sec-formal-command-grammar}) and remains fully in scope ---
  the two commands share a letter-adjacent name but not an implementation,
  and this bullet excludes only the sticky multi-character mode.
\item
  \textbf{Folds} --- \texttt{zf}, \texttt{zo}, \texttt{zc}, and fold-aware motions.
\item
  \textbf{Additional text objects} --- sentence objects (\texttt{is}/\texttt{as}), and any
  syntax- or treesitter-aware object beyond the fixed set in Section~\ref{sec-text-objects}'s
  table.
\item
  \textbf{The Ex-command language} --- beyond the two named exceptions below, full
  Ex-scripting, \texttt{:let}, arbitrary ranges beyond whole-buffer/global-matched
  lines, and anything requiring Vimscript evaluation is out of scope.
  \textbf{Two explicit exceptions}, specified in full in Section~\ref{sec-global-commands-and-substitute}, exist
  because real usage showed the original single-command carve-out (macro
  replay only) was too narrow to cover how \texttt{:g} and \texttt{:s} actually get used
  together in practice: the global-command family
  (\texttt{:g}/\texttt{:global}, \texttt{:v}/\texttt{:g!}/\texttt{:vglobal}) combined with a small fixed set of
  per-line actions (\texttt{d}, \texttt{m}, \texttt{t}/\texttt{co}, \texttt{normal{[}!{]}}), and the substitute
  command (\texttt{:s}, \texttt{:\%s}) with a bounded pattern/replacement grammar that
  explicitly excludes the \texttt{\textbackslash{}=} expression register. These are two named
  additions with their own specification, not a reopening of the
  Ex-command family generally.
\item
  \textbf{Plugins and Vimscript execution} --- this project teaches Vim's built-in
  grammar; it is not a Vimscript runtime. This is also why the \texttt{\textbackslash{}=}
  expression register is excluded from the \texttt{:substitute} grammar in
  Section~\ref{sec-global-commands-and-substitute} even though \texttt{:substitute} itself is in scope: \texttt{\textbackslash{}=} evaluates
  an arbitrary expression (\texttt{printf(...)}, \texttt{strftime(...)}, \texttt{line(\textquotesingle{}.\textquotesingle{})}) as
  part of the replacement, and is exactly the Vimscript-execution surface
  this exclusion is naming, just reachable through a single command's
  syntax rather than through a \texttt{:let}-style statement.
\item
  \textbf{Treesitter-based or language-aware text objects} --- anything requiring
  actual syntax parsing of the buffer's contents as a programming language,
  as opposed to the character-class rules in Section~\ref{sec-word-classification}.
\item
  \textbf{Bash/readline ``vi mode'' (\texttt{set\ -o\ vi}) is not Vim and is not a coverage
  target at all}, named here only to close off a plausible confusion: it
  is a different program's emulation of Vim's modal editing, with real
  behavioral differences (no multi-line editing, a different and smaller
  text-object set) that would violate the Section~\ref{sec-fidelity-principle} Fidelity Principle if
  treated as interchangeable with the actual Vim behavior this document
  specifies.
\end{itemize}

Each of the still-excluded items above is a plausible ``Phase 9'' for a
version of this project that outgrows its original teaching scope, not a
rejection of the idea --- they are listed here specifically so that scope
creep during Phases 0--8 has a named list to be checked against and deferred
to, rather than being negotiated fresh each time it comes up.

\begin{center}\rule{0.5\linewidth}{0.5pt}\end{center}

\hypertarget{sec-global-commands-and-substitute}{%
\section{Global Commands and Substitute}\label{sec-global-commands-and-substitute}}

Section~\ref{sec-future-vim-coverage-explicitly-out-of-scope-for-v1} named two exceptions to the Ex-command exclusion. This section
specifies both in full, at the same level of rigor as Section~\ref{sec-formal-command-grammar}, 4, and 5,
since real usage showed both are load-bearing rather than peripheral.

\hypertarget{sec-the-global-command-family}{%
\subsection{The Global Command Family}\label{sec-the-global-command-family}}

\begin{Verbatim}[breaklines=true, breakanywhere=true, fontsize=\small]
<global>        ::= <global-normal> | <global-inverse>
<global-normal> ::= (":g" | ":global") "/" <pattern> "/" <line-action>
<global-inverse>::= (":v" | ":g!" | ":vglobal") "/" <pattern> "/" <line-action>
<line-action>   ::= "d"                            ; delete every matching line
                   | "m" <dest-line>                ; move every matching line to dest-line
                   | ("t" | "co") <dest-line>       ; copy every matching line to after dest-line
                   | "normal!"? <keys>              ; run a keystroke sequence at each match
                   | <substitute>                   ; see 12.2 — ":g/pat/s/.../.../ " form
\end{Verbatim}

\texttt{\textless{}pattern\textgreater{}} is the same pattern grammar as Section~\ref{sec-substitute}'s search side. The
semantics that matter for correctness, and are easy to get wrong:

\begin{itemize}
\tightlist
\item
  \texttt{:g} first scans the \textbf{entire buffer once}, marking every line that
  matches, and only then executes \texttt{\textless{}line-action\textgreater{}} on each marked line in
  order --- it does not re-scan after each action. This is why
  \texttt{:g/pattern/normal\ \$3X} behaves correctly even when the action changes
  line lengths or content: the set of target lines was fixed before any
  action ran.
\item
  \texttt{:g/\^{}/m0} (move every line to line 0, in buffer order) is a real,
  idiomatic reverse-the-buffer idiom precisely because of the
  scan-then-act ordering above: each line, moved to position 0 in turn,
  ends up in reverse order relative to the original.
\item
  \texttt{normal!} (with the bang) suppresses the user's own custom mappings
  during the per-line replay, using only the grammar this document
  specifies; \texttt{normal} without the bang is excluded from this grammar, since
  it would depend on a remapping layer this project does not implement.
\item
  A \texttt{\textless{}keys\textgreater{}} sequence following \texttt{normal!} may itself invoke a macro
  register (\texttt{@a}) --- this is the original narrow exception from the prior
  revision of Section~\ref{sec-future-vim-coverage-explicitly-out-of-scope-for-v1} --- but it is no longer the \emph{only} supported
  \texttt{\textless{}line-action\textgreater{}}, since \texttt{d}, \texttt{m}, \texttt{t}/\texttt{co}, and direct keystrokes without
  a macro are equally common and equally in scope.
\item
  \textbf{Nested \texttt{:g}/\texttt{:v} is illegal, matching real Vim.} Neither \texttt{\textless{}keys\textgreater{}}
  (following \texttt{normal!}) nor \texttt{\textless{}substitute\textgreater{}} may itself contain a \texttt{:g},
  \texttt{:global}, \texttt{:v}, \texttt{:g!}, or \texttt{:vglobal} invocation. Real Vim rejects this
  with \texttt{E147:\ Cannot\ do\ :global\ recursive}, and this grammar rejects it the
  same way: a parse-time error, not a silently-ignored or partially-applied
  command. This matters specifically because \texttt{\textless{}line-action\textgreater{}} already
  permits \texttt{normal!} with an arbitrary \texttt{\textless{}keys\textgreater{}} sequence, which is exactly
  general enough that an implementer could accidentally allow a second
  \texttt{:g} to slip through inside it if this exclusion weren't stated
  separately from the rest of the grammar.
\end{itemize}

\hypertarget{sec-substitute}{%
\subsection{Substitute}\label{sec-substitute}}

\begin{Verbatim}[breaklines=true, breakanywhere=true, fontsize=\small]
<substitute>  ::= <range>? "s" <sep> <pattern> <sep> <replacement> <sep> <flags>?
<range>       ::= "%"                    ; whole buffer — the primary case
                 | <line-num> "," <line-num>
<sep>         ::= "/" | "#" | "!" | ...  ; any single non-alphanumeric char, consistent within one command
<flags>       ::= ("g" | "i")+           ; g: all matches per line, not just first; i: case-insensitive
\end{Verbatim}

\textbf{No \texttt{\textless{}range\textgreater{}} means the current line only, not the whole buffer.} The \texttt{?}
on \texttt{\textless{}range\textgreater{}} in the grammar above is easy to misread as ``range is optional,
so omitting it is the same as the widest range'' --- this is backwards. A bare
\texttt{:s/foo/bar/} with no range prefix applies to the cursor's current line
only; \texttt{\%} must be given explicitly to mean the whole buffer, exactly as
\texttt{\textless{}range\textgreater{}} requires it as one of two named alternatives rather than as the
grammar's default. This is worth stating outside the grammar because it is
the single most common mistake a Vim beginner makes with \texttt{:s}, and a
teaching tool that got it wrong would be teaching the mistake rather than
correcting it.

\textbf{Pattern grammar (search side), in scope:}

\begingroup\small
\begin{longtable}[]{@{}
  >{\raggedright\arraybackslash}p{(\columnwidth - 6\tabcolsep) * \real{0.2500}}
  >{\raggedright\arraybackslash}p{(\columnwidth - 6\tabcolsep) * \real{0.2500}}
  >{\raggedright\arraybackslash}p{(\columnwidth - 6\tabcolsep) * \real{0.2500}}
  >{\raggedright\arraybackslash}p{(\columnwidth - 6\tabcolsep) * \real{0.2500}}@{}}
\toprule\noalign{}
\begin{minipage}[b]{\linewidth}\raggedright
Feature
\end{minipage} & \begin{minipage}[b]{\linewidth}\raggedright
Magic mode
\end{minipage} & \begin{minipage}[b]{\linewidth}\raggedright
Very-magic mode (\texttt{\textbackslash{}v})
\end{minipage} & \begin{minipage}[b]{\linewidth}\raggedright
Example from real usage
\end{minipage} \\
\midrule\noalign{}
\endhead
\bottomrule\noalign{}
\endlastfoot
literal text & as written & as written & \\
any character & \texttt{.} & \texttt{.} & \\
quantifiers & \texttt{\textbackslash{}+} \texttt{\textbackslash{}*} \texttt{\textbackslash{}\{n,m\}} & \texttt{+} \texttt{*} \texttt{\{n,m\}} & \texttt{\textbackslash{}d\textbackslash{}\{1,3\}} \\
character classes & \texttt{\textbackslash{}d} \texttt{\textbackslash{}w} \texttt{\textbackslash{}s} \texttt{\textbackslash{}u} & \texttt{\textbackslash{}d} \texttt{\textbackslash{}w} \texttt{\textbackslash{}s} \texttt{\textbackslash{}u} & \texttt{\textbackslash{}d\textbackslash{}\{1,3\}\textbackslash{}.\textbackslash{}d\textbackslash{}\{1,3\}...} \\
anchors & \texttt{\^{}} \texttt{\$} & \texttt{\^{}} \texttt{\$} & \\
grouping & \texttt{\textbackslash{}(...\textbackslash{})} & \texttt{(...)} & \texttt{\textbackslash{}v\^{}(a\textbackslash{}\textbar{}g)} \\
alternation & \texttt{\textbackslash{}\textbar{}} (needs \texttt{\textbackslash{}}) & \texttt{\textbackslash{}\textbar{}} & \texttt{\textbackslash{}v\^{}(a\textbackslash{}\textbar{}g)} \\
match-start reset & \texttt{\textbackslash{}zs} & \texttt{\textbackslash{}zs} & \texttt{\textbackslash{}v(\^{}\textbackslash{}\textbar{}{[}.!?{]}\ )\textbackslash{}zs\textbackslash{}w} \\
\end{longtable}
\endgroup

Both magic levels are in scope, not just one, because real usage mixes them
(default-magic \texttt{\textbackslash{}d\textbackslash{}\{1,3\}} alongside very-magic \texttt{\textbackslash{}v\^{}(a\textbar{}g)}), and a learner
who only sees one style in-game would be underprepared for the other.

\textbf{Replacement grammar, in scope:} literal text, whole-match backreference
\texttt{\&}, numbered group backreferences \texttt{\textbackslash{}0}--\texttt{\textbackslash{}9}, and the case-modifier escapes
\texttt{\textbackslash{}u} \texttt{\textbackslash{}U} \texttt{\textbackslash{}l} \texttt{\textbackslash{}L} \texttt{\textbackslash{}e} \texttt{\textbackslash{}E} (next-char, run-until-\texttt{\textbackslash{}E}, and end-of-run
case conversion --- this is exactly the mechanism behind the small-caps and
sentence-capitalization idioms real usage showed: \Verb[breaklines=true,breakanywhere=true]#%s/\u\+/\L&/g# and
\Verb[breaklines=true,breakanywhere=true]#%s/\v(^|[.!?] )\zs\w/\u&/g#).

\textbf{Explicitly excluded from the replacement grammar:} the \texttt{\textbackslash{}=} expression
register (\texttt{\textbackslash{}=printf(...)}, \texttt{\textbackslash{}=strftime(...)}, \texttt{\textbackslash{}=line(\textquotesingle{}.\textquotesingle{})}, or any other
expression). This is a single, deliberate cut: \texttt{\textbackslash{}=} is the one point in
\texttt{:substitute}'s own syntax where arbitrary Vimscript evaluation becomes
reachable, and excluding it removes that whole surface at once without
narrowing anything else in this section's grammar. A level or macro that
depends on \texttt{\textbackslash{}=} is out of scope by construction, the same as any other
Vimscript dependency named in Section~\ref{sec-future-vim-coverage-explicitly-out-of-scope-for-v1}.

\begin{center}\rule{0.5\linewidth}{0.5pt}\end{center}

\hypertarget{sec-user-generated-and-editable-levels}{%
\section{User-Generated and Editable Levels}\label{sec-user-generated-and-editable-levels}}

The curriculum levels in Section~\ref{sec-level-data-format} are hand-authored and pedagogically
ordered. This section specifies a second, complementary level type: an
arbitrary real editing task --- reformat this data, fix this repeated typo,
rename this pattern everywhere --- imported and played as a level in its own
right, scored the same way. This is the design's actual payoff: once a task
can be expressed as a \texttt{before}/\texttt{after} buffer pair, the game doesn't need to
know or care what the task was \emph{for}.

\hypertarget{sec-level-import-model}{%
\subsection{Level Import Model}\label{sec-level-import-model}}

Import needs only what Section~\ref{sec-level-data-format}'s \texttt{goal} field already accepts: a starting
buffer and a target. Concretely:

\begin{itemize}
\tightlist
\item
  \textbf{Direct pair.} The author supplies \texttt{initialBuffer} and a target buffer
  directly; the target becomes a \texttt{bufferEquals} goal exactly as in Section~\ref{sec-level-data-format}.
\item
  \textbf{Diff-based.} For large real files, pasting two full copies is unwieldy.
  The author instead supplies \texttt{initialBuffer} and a unified diff; the target
  buffer is computed once at import time by applying the diff, and only the
  resulting target is stored in the level definition. The diff itself is an
  authoring convenience --- it is not part of the level's runtime
  representation, and nothing downstream (parser, oracle, leaderboard) ever
  sees it.
\end{itemize}

No new parser feature is required for either path; this section is an
authoring workflow layered on Section~\ref{sec-level-data-format}'s existing schema, not a new goal
type.

\hypertarget{sec-author-plays-first-par-seeding}{%
\subsection{Author-Plays-First Par Seeding}\label{sec-author-plays-first-par-seeding}}

Section~\ref{sec-level-data-format}'s \texttt{par} field assumed a level designer who simply knows the
intended solution length. That doesn't generalize to arbitrary imported
tasks, since computing a provably minimal keystroke sequence over the full
grammar (counts, registers, macros, the \texttt{:g} exception above) is not
tractable in general. Par is therefore seeded operationally rather than
computed: \textbf{the author must solve the level themselves once at import
time}, and their own keystroke count becomes the level's starting par. This
guarantees, as a side effect of how the level is created, that at least one
real solution exists --- a level that turns out to be unsolvable as authored
is caught at the moment of authoring, not after a player gets stuck.

\hypertarget{sec-crowd-sourced-par-and-verified-replays}{%
\subsection{Crowd-Sourced Par and Verified Replays}\label{sec-crowd-sourced-par-and-verified-replays}}

Par for an imported level is a discovered record, not a computed constant ---
the same posture Pipeline Golf takes toward transformation quality elsewhere
in this line of work: no claimed optimum, only a current best that stands
until beaten.

This only works if submissions are trustworthy, which means a submitted
score can never be a bare number:

\begin{itemize}
\tightlist
\item
  Every submission stores the \textbf{actual keystroke sequence} that achieved
  it, not just its length.
\item
  A submission is accepted only if replaying that exact sequence through the
  deterministic, headless parser (Sections~\ref{sec-mode-state-machine}--\ref{sec-text-objects}) reaches the level's declared
  goal state in precisely the claimed number of keystrokes. This reuses the
  same replay mechanism Section~\ref{sec-differential-testing-against-real-vim} already requires for oracle testing ---
  both are ``replay these keys against this parser and check the result,''
  differing only in which state they check against (real Neovim's output vs.
  a level's declared goal).
\end{itemize}

\textbf{What counts as one keystroke is not obvious, and a par system built on an
ambiguous definition isn't a par system.} Two implementations that
disagree on whether \texttt{f\ a} is one keystroke or two, or whether invoking a
macro costs 1 keystroke or the N keystrokes it replays, will produce
different par values for the \emph{same} recorded solution --- which defeats the
purpose of a shared, cross-verifiable leaderboard. The rule is: \textbf{every
distinct key-press event that reaches the parser as a discrete input counts
as exactly one keystroke}, with no event ever counted as zero, as a
fraction, or as a multiple of what was physically pressed. The table below
resolves every case this document's grammar makes possible; the general
rule above is what a new case should be checked against if one is missed.

\begingroup\small
\begin{longtable}[]{@{}
  >{\raggedright\arraybackslash}p{(\columnwidth - 4\tabcolsep) * \real{0.3333}}
  >{\raggedright\arraybackslash}p{(\columnwidth - 4\tabcolsep) * \real{0.3333}}
  >{\raggedright\arraybackslash}p{(\columnwidth - 4\tabcolsep) * \real{0.3333}}@{}}
\toprule\noalign{}
\begin{minipage}[b]{\linewidth}\raggedright
Input
\end{minipage} & \begin{minipage}[b]{\linewidth}\raggedright
Keystroke count
\end{minipage} & \begin{minipage}[b]{\linewidth}\raggedright
Reasoning
\end{minipage} \\
\midrule\noalign{}
\endhead
\bottomrule\noalign{}
\endlastfoot
A single command key (\texttt{d}, \texttt{w}, \texttt{x}, \texttt{\textasciitilde{}}, \texttt{\$}, \texttt{\%}, \texttt{\{}, \texttt{.}) & 1 & One key-press event \\
\texttt{\textless{}Esc\textgreater{}} & 1 & One key-press event, same as any other key \\
A control-chord (\texttt{Ctrl-r}, \texttt{Ctrl-v}, \texttt{Ctrl-w}, \texttt{Ctrl-{[}}) & 1 & The chord is a single input event at the keyboard level; it is not counted per physical key held \\
Each digit of a count prefix (\texttt{12dw}) & 1 per digit & \texttt{12dw} = \texttt{1},\texttt{2},\texttt{d},\texttt{w} = \textbf{4}, not 3 --- a two-digit count costs exactly as much as physically typing two digits, which is what makes a needlessly large count an honest cost rather than a free multiplier \\
\texttt{f\{char\}}, \texttt{F\{char\}}, \texttt{t\{char\}}, \texttt{T\{char\}}, \texttt{r\{char\}} & 2 & The command key and the target character are separate key-press events \\
Each character of a search pattern, including the leading \texttt{/} or \texttt{?} and the terminating \texttt{\textless{}CR\textgreater{}} & 1 per character & \texttt{/foo\textless{}CR\textgreater{}} = \texttt{/},\texttt{f},\texttt{o},\texttt{o},\texttt{\textless{}CR\textgreater{}} = \textbf{5} \\
Each character of an Ex command line, including the leading \texttt{:} and the terminating \texttt{\textless{}CR\textgreater{}} & 1 per character & \texttt{:\%s/a/b/g\textless{}CR\textgreater{}} = \textbf{12} --- this is what makes the par system fairly weigh ``type a substitute command'' against ``edit by hand,'' rather than giving \texttt{:s} an unearned discount \\
Each typed character of inserted text, plus the mode-entering key and the terminating \texttt{\textless{}Esc\textgreater{}} & 1 per character/key & \texttt{ihi\textless{}Esc\textgreater{}} = \texttt{i},\texttt{h},\texttt{i},\texttt{\textless{}Esc\textgreater{}} = \textbf{4} \\
A text object or motion spelled with multiple characters (\texttt{iw}, \texttt{aw}, \texttt{ge}) & 1 per character & \texttt{diw} = \texttt{d},\texttt{i},\texttt{w} = \textbf{3} \\
Macro recording (\texttt{qa\ ...\ q}) & 1 for \texttt{q}, 1 for the register letter, 1 for each keystroke performed while recording, 1 for the terminating \texttt{q} & Recording has no discount: it costs exactly what was typed, including the \texttt{q}/register/\texttt{q} bracketing \\
Macro invocation (\texttt{@a}) & 2 & \texttt{@} and the register letter, regardless of how many keystrokes the macro replays internally --- this is the entire reason macros are efficient in this accounting, not an oversight to patch \\
Repeated macro invocation (\texttt{10@a}) & 1 per digit, plus 2 & \texttt{10@a} = \texttt{1},\texttt{0},\texttt{@},\texttt{a} = \textbf{4}, however many keystrokes the macro itself replays each of the ten times \\
\texttt{.} (dot-repeat) & 1 & Regardless of the size or complexity of the command being repeated --- this is symmetric with macro invocation: both are cases where one keystroke stands for an arbitrarily larger replayed effect \\
\end{longtable}
\endgroup

A worked example ties this to Section~\ref{sec-the-global-command-family}'s own idiom: \texttt{:g/\^{}\textbackslash{}d\textbackslash{}+\$/d} ---
every character including \texttt{:}, \texttt{g}, \texttt{/}, \texttt{\^{}}, \texttt{\textbackslash{}}, \texttt{d}, \texttt{+}, \texttt{\$}, \texttt{/}, \texttt{d},
and the terminating \texttt{\textless{}CR\textgreater{}} counts individually, for a total of \textbf{11}
keystrokes, regardless of how many lines in the buffer actually match and
get deleted. This is the same principle as the Ex-command row above,
applied to \texttt{:g} rather than \texttt{:s}: the cost is what was typed, not what the
command's effect turned out to be.

\begin{itemize}
\tightlist
\item
  Leaderboards are naturally shared, per-level data: \texttt{window.storage} with
  \texttt{shared:\ true}, keyed by level id, storing an ordered list of
  \texttt{\{keystrokes:\ string,\ playerLabel:\ string\}} entries rather than scores
  alone, so any entry remains independently re-verifiable by replaying it ---
  including, as a spot check, by running it through the Section~\ref{sec-differential-testing-against-real-vim} oracle
  as well as the level's own goal check.
\end{itemize}

\hypertarget{sec-repetition-detection-and-macro-friendliness}{%
\subsection{Repetition Detection and Macro-Friendliness}\label{sec-repetition-detection-and-macro-friendliness}}

Real editing tasks that are worth turning into levels are disproportionately
the ones with repeated structure --- the same shape of edit recurring at many
sites --- which is exactly the condition under which a short macro dominates
any manual approach. This is the actual reason this pivot makes macros the
star mechanic rather than one system among several.

At import time, a lightweight structural diff between \texttt{initialBuffer} and
the target can detect approximately-repeated edit patterns and annotate the
level with an estimated repetition count. This is metadata for hinting and
level browsing (``macro-friendly,'' tagged with an approximate multiplicity),
not a gameplay gate --- \texttt{allowedCommands} and the lock system from Section~\ref{sec-pedagogical-layer}
still govern what's actually available in a given level. Detection only
decides what the level is labeled, not what it permits.

\hypertarget{sec-placement-in-the-build-plan}{%
\subsection{Placement in the Build Plan}\label{sec-placement-in-the-build-plan}}

This section depends on the full grammar being solid --- specifically macros
(Phase 6) and, for large imported files, the undo-tree work (Phase 7) --- and
benefits from Section~\ref{sec-differential-testing-against-real-vim}'s oracle harness already existing, since the same
replay-and-verify mechanism does double duty as the leaderboard's anti-cheat
check. It belongs after the curriculum phases as its own \textbf{Phase 8}, so
that level-authoring tooling never becomes a dependency the core teaching
path in Phases 0--7 has to wait on.

\end{document}
